% !PRINTED 340
\weberpage{340}
%<<CONT

3. \gesperrt{If $\varDelta$ is positive, $\sqrt{\varDelta}$ shall likewise be
positive, and if $\varDelta$ is negative, $-i\sqrt{\varDelta}$ shall be positive, or
$\sqrt{\varDelta}$ positive imaginary.}

From (1), (2) there follows further:

4. \gesperrt{The fundamental number of a quadratic field is
always $\equiv 0$ or $\equiv 1$ with respect to the modulus 4, and has, apart from 4,
no quadratic divisor. It is therefore a discriminant,
and indeed a fundamental discriminant.}

\textbf{\S\ 91. \gesperrt{Ideals and Forms in Quadratic Fields.}}

If $\mathfrak{a}$ is an ideal of the quadratic field $\Omega$, $\alpha_1, \alpha_2$ a
positive basis, and
\weq{1}{\lambda=\alpha_1 t_1+\alpha_2 t_2}
a basis form of $\mathfrak{a}$, then
\weq{2}{N(\lambda)=(\alpha_1 t_1+\alpha_2 t_2)(\alpha_1' t_1+\alpha_2' t_2)=N(\mathfrak{a})\,T,}
and $N(\mathfrak{a})$ is the greatest common divisor of the three integral
rational numbers
\[
\alpha_1\alpha_1', \quad \alpha_1\alpha_2'+\alpha_2\alpha_1', \quad \alpha_2\alpha_2'.
\]
If we therefore set
\weq{3}{\begin{aligned}
\alpha_1\alpha_1' &= a\,N(\mathfrak{a}),\\
\alpha_1\alpha_2'+\alpha_2\alpha_1' &= b\,N(\mathfrak{a}),\\
\alpha_2\alpha_2' &= c\,N(\mathfrak{a}),
\end{aligned}}
then
\weq{4}{T=a t_1^2+b t_1 t_2+c t_2^2}
is the form belonging to the basis $\alpha_1, \alpha_2$. We denote it, when
the notation for the variables is of no consequence, by
\weq{5}{\varphi=(a, b, c).}
It is a primitive quadratic form whose discriminant
\weq{6}{\varDelta=b^2-4ac}
is equal to the fundamental number of the field $\Omega$ (\S\ 87). In order to form
the basis $\alpha_1, \alpha_2$ according to Bd.\,II, \S\ 163, one seeks first the least
natural number $a_{11}$ divisible by $\mathfrak{a}$, and then the least
natural number $a_{12}$ for which a rational number $a_{12}$ satisfying the condition % [sic] German source prints a_{12} in both occurrences.

\weq{7}{a_{12}+a_{22}\theta\equiv 0\pmod{\mathfrak{a}}}
can be determined. Since the congruence (7) can be satisfied for
$a_{22}=a_{11}$, namely by $a_{12}=0$, then
%CONT>>
